% !TEX root = ../../main.tex
\subsection{场景创新}

现有外卖平台通常已经具备用户与骑手之间的即时通信能力，但其核心目标主要是提升配送效率，而非系统性解决通信过程中的隐私保护与证据可信问题。在实际外卖场景中，用户与骑手之间的通信往往伴随订单号、配送地址、联系方式、备注信息和聊天内容等多类敏感信息。如果系统缺少精细化的权限边界，骑手端或管理员端可能接触到超出配送所需范围的信息；如果系统缺少可信审计机制，平台在处理纠纷、骚扰投诉或恶意争议时又难以证明通信记录是否保持原始状态。

FoodShield 的场景创新在于：将外卖订单通信从普通业务聊天场景提升为一个兼具隐私保护、身份认证和审计治理需求的安全场景。系统以订单为基本单位组织通信、认证和审计流程，使订单号不仅是业务编号，同时也是通信房间、认证凭证、日志记录和条件溯源的统一索引。用户与骑手的日常通信在匿名保护下完成，异常场景下则通过订单号触发受控溯源流程，从而避免了传统方案中“直接实名导致隐私暴露”和“完全匿名导致无法治理”两类极端问题。

此外，本作品选择外卖平台这一高频、真实、贴近日常生活的应用场景，能够较直观地展示信息安全技术在生活服务平台中的实际价值。相比单纯的算法演示或通用匿名聊天系统，本作品更强调安全机制与具体业务流程之间的结合，使系统创新具有更强的可解释性和展示效果。
